\SourcePage{686}{689}
\section{Collisions of identical particles}

The collision of two identical particles requires separate consideration. In quantum mechanics, particle identity produces a peculiar exchange interaction between them, which has a substantial effect on scattering (N.~Mott, 1930)\footnote{The direct spin--orbit interaction is still not considered here.}.

For a two-particle system, the orbital wavefunction must be symmetric or antisymmetric under interchange of the particles, according as their combined spin is even or odd (see \S~62). Consequently, the scattering wavefunction obtained by solving the ordinary Schrödinger equation must likewise be symmetrized or antisymmetrized in the particle labels. Interchanging the particles amounts to reversing the radius vector that joins them. In the frame where the centre of mass is at rest, this leaves $r$ unchanged while replacing the angle $\theta$ by $\pi-\theta$ (and hence sends $z=r\cos\theta$ to $-z$). Thus, in place of the asymptotic expression (123.3) for the wavefunction, we must write

\begin{equation}
 \psi=e^{ikz}\pm e^{-ikz}+\frac1r e^{ikr}[f(\theta)\pm(\pi-\theta)].
 \tag{137.1}
\end{equation}
Because the particles are identical, one cannot say which is scattered and which is the scatterer. In the centre-of-mass frame there are two identical incident waves propagating towards each other, $e^{ikz}$ and $e^{-ikz}$. The outgoing spherical wave in (137.1) accounts for the scattering of both particles, and the flux calculated from it gives the probability that either particle is scattered into a given solid-angle element $d o$. The scattering cross-section is the ratio of this flux to the flux density in each incident

\SourcePage{687}{690}
plane wave; as before, it is therefore the squared modulus of the coefficient of $e^{ikr}/r$ in (137.1).

If the total spin of the colliding particles is even, the cross-section is
\begin{equation}
 d\sigma_s=|f(\theta)+f(\pi-\theta)|^2d o,
 \tag{137.2}
\end{equation}
while for odd total spin it is
\begin{equation}
 d\sigma_a=|f(\theta)-f(\pi-\theta)|^2d o.
 \tag{137.3}
\end{equation}
The appearance of the interference term $f(\theta)f^*(\pi-\theta)+f^*(\theta)f(\pi-\theta)$ is characteristic of the exchange interaction. If the particles were distinguishable, as in ordinary classical mechanics, the probability for either of them to be scattered into $d o$ would simply be the sum of the probabilities for one to be deflected through $\theta$ and the oppositely moving one through $\pi-\theta$. The cross-section would then be
\[
 \{|f(\theta)|^2+|f(\pi-\theta)|^2\}d o.
\]

In the low-velocity limit, the scattering amplitude, provided the interaction decreases sufficiently rapidly with distance, tends to an angle-independent constant (\S~132). Equation (137.3) then shows that $d\sigma_a$ vanishes: only particles with even total spin scatter from one another.

Equations (137.2) and (137.3) assume a definite value for the combined spin of the colliding particles. If the particles are not in definite spin states, the cross-section must instead be found by averaging over all spin states, taking them to be equally probable.


It was shown in \S~62 that among the $(2s+1)^2$ spin states of two particles of spin $s$, $s(2s+1)$ states have even total spin and $(s+1)(2s+1)$ have odd total spin if $s$ is half-integral, with the assignments reversed if $s$ is integral. First let $s$ be half-integral. The probabilities of even and odd $S$ are then $s/(2s+1)$ and $(s+1)/(2s+1)$, respectively. Hence
\begin{equation}
 d\sigma=\frac{s}{2s+1}d\sigma_s+\frac{s+1}{2s+1}d\sigma_a.
 \tag{137.4}
\end{equation}

\SourcePage{688}{691}
Substitution of (137.2) and (137.3) gives
\begin{equation}
 d\sigma\left\{|f(\theta)|^2+|f(\pi-\theta)|^2
 -\frac1{2s+1}[f(\theta)f(\pi-\theta)^*+f(\theta)^*f(\pi-\theta)]\right\}d o.
 \tag{137.5}
\end{equation}
For integral $s$, similarly,
\begin{equation}
 d\sigma=\left\{|f(\theta)|^2+|f(\pi-\theta)|^2
 +\frac1{2s+1}[f(\theta)f(\pi-\theta)^*+f(\theta)^*f(\pi-\theta)]\right\}d o.
 \tag{137.6}
\end{equation}

As an example, consider two electrons interacting through the Coulomb law $U=e^2/r$. Substitution of (135.9) in (137.5) with $s=1/2$ gives, after a simple calculation in ordinary units,
\begin{align}
 d\sigma={}&\left(\frac{e^2}{m_0v^2}\right)^2
 \left[\frac1{\sin^4(\theta/2)}+\frac1{\cos^4(\theta/2)}\right.\notag\\
 &\left.-\frac1{\sin^2(\theta/2)\cos^2(\theta/2)}
 \cos\left(\frac{e^2}{\hbar v}\ln\tan^2\frac\theta2\right)\right]d o
 \tag{137.7}
\end{align}
(we have introduced the electron mass $m_0$ in place of the reduced mass $m=m_0/2$). This formula simplifies appreciably when the speed is so high that $e^2\ll v\hbar$, precisely the condition for perturbation theory to apply to the Coulomb field. The cosine in the third term can then be replaced by unity, giving
\begin{equation}
 d\sigma=\left(\frac{2e^2}{m_0v^2}\right)^2\frac{4-3\sin^2\theta}{\sin^4\theta}\,d o.
 \tag{137.8}
\end{equation}
The opposite limit, $e^2\gg v\hbar$, corresponds to the classical limit (see the end of \S~127). In (137.7) this limit is reached in a distinctive manner. For $e^2\gg v\hbar$, the cosine in the third term oscillates rapidly. At any prescribed $\theta$, (137.7) generally gives a cross-section appreciably different from Rutherford's. Averaging over even a small interval of $\theta$, however, removes the oscillatory term and gives the classical formula.

All these formulas refer to the centre-of-mass frame. Passage to a frame in which one particle was at rest before the collision is accomplished,

\SourcePage{689}{692}
according to (123.2), simply by replacing $\theta$ by $2\vartheta$. Thus (137.7) gives for electron collisions
\begin{align}
 d\sigma={}&\left(\frac{2e^2}{m_0v^2}\right)^2
 \left[\frac1{\sin^4\vartheta}+\frac1{\cos^4\vartheta}
 -\frac1{\sin^2\vartheta\cos^2\vartheta}\right.\notag\\
 &\left.\times\cos\left(\frac{e^2}{\hbar v}\ln\tan^2\vartheta\right)\right]
 \cos\vartheta\,d o,
 \tag{137.9}
\end{align}
where $d o$ is the solid-angle element in the new frame. When $\theta$ is replaced by $2\vartheta$, $d o$ must be replaced by $4\cos\vartheta\,d o$, since $\sin\theta\,d\theta\,d\varphi=4\cos\vartheta\sin\vartheta\,d\vartheta\,d\varphi$.

\subsection*{Problem}

Determine the scattering cross-section of two identical spin-$1/2$ particles having prescribed mean spins $\bar{\mathbf s}_1$ and $\bar{\mathbf s}_2$.

\SolutionHeading The dependence of the cross-section on the particle polarizations must occur through a term proportional to the scalar $\bar{\mathbf s}_1\bar{\mathbf s}_2$. Seek $d\sigma$ in the form $a+b\bar{\mathbf s}_1\bar{\mathbf s}_2$. For unpolarized particles, $\bar{\mathbf s}_1=\bar{\mathbf s}_2=0$, the second term is absent and (137.4) gives $d\sigma=a=(d\sigma_s+3d\sigma_a)/4$. If both particles are fully polarized in the same direction, $\bar{\mathbf s}_1\bar{\mathbf s}_2=1/4$, the system is certainly in a state with $S=1$; hence $d\sigma=a+b/4=d\sigma_a$. Solving these two equations for $a$ and $b$, we find
\[
 d\sigma=\frac14(d\sigma_s+3d\sigma_a)+(d\sigma_a-d\sigma_s)\bar{\mathbf s}_1\bar{\mathbf s}_2.
\]
